\documentclass[12pt]{article}
\usepackage{amsmath, amssymb}
\usepackage{geometry}
\usepackage{hyperref}
\usepackage{parskip}
\usepackage{microtype}
\usepackage{booktabs}
\usepackage{xcolor}
\geometry{margin=1.3in}

\definecolor{gold}{RGB}{180,140,30}

\title{\textbf{Quark Masses in the Unified Mechanics Framework:\\
The $\varphi^4$ Cycle Pattern, Isospin Splitting,\\
and the Accumulated Braiding Floor}}
\author{Joseph Shields}
\date{2026}

\begin{document}
\maketitle

\begin{abstract}
The Unified Mechanics axiom $\varphi^2 = \varphi + 1$ predicts the strong coupling
$\alpha_s(M_Z) = r^2(1+r-r^2) \approx 0.1159$ ($-1.71\%$,
$0.58\,\varepsilon_{\text{floor}}$). The quark masses are organised by
a $\varphi^4$ colour-cycle pattern: each quark sits within $1.6$ cycles of
the nearest integer multiple of 4 in $\log_\varphi(m_q/m_e)$ space.
When residuals are normalised by the accumulated braiding floor
$n\,\varepsilon_{\text{floor}}$ (where $n$ is the cycle depth), $c$, $b$, and
$t$ fall within $1.5\,\varepsilon_{\text{floor}}$. The $u$--$d$ isospin
splitting is derived from the electroweak channel structure:
$m_d/m_u = \sqrt{W_B/W_M} = 1/\!\sin\theta_W \approx 2.115$ (PDG 2.162,
$-2.2\%$, $0.74\,\varepsilon_{\text{floor}}$). With this correction the
individual $u$ and $d$ masses each fall within $1.4\,\varepsilon_{\text{floor}}$.
The strange quark (at a distinct 4-cycle depth from charm) is treated via a
boundary-channel coupling factor $\sqrt{W_B}$, giving $m_s = \varphi^{12}m_e\sqrt{W_B}
= 107.5\,\text{MeV}$ ($+15.1\%$, $0.43\,n\,\varepsilon_{\text{floor}}$, $n = 12$),
consistent with the accumulated floor. The PDG $1\sigma$ uncertainty of
$\pm 11\,\text{MeV}$ spans the UM prediction; derivation of the $\sqrt{W_B}$
factor from $G_2$ first principles is the next calculation.
\end{abstract}

%% ────────────────────────────────────────────────────────────
\section{Foundations}
%% ────────────────────────────────────────────────────────────

The axiom and its immediate consequences (Paper~1):
\begin{equation}
  \varphi^2 = \varphi + 1
  \quad\Rightarrow\quad
  \varphi = \tfrac{1+\sqrt{5}}{2},\qquad
  r = \tfrac{1}{2\varphi} \approx 0.30902.
  \label{eq:axiom}
\end{equation}

Channel weights and braiding floor:
\begin{align}
  W_L &= (1-r)^2 \approx 0.4775, &
  W_B &= 2r(1-r) \approx 0.4271, &
  W_M &= r^2 \approx 0.0955, \\
  \varepsilon_{\text{floor}} &= r^3 \approx 0.02951.
\end{align}

The strong coupling from the colour-channel weight structure (Paper~5):
\begin{equation}
  \alpha_s(M_Z) = r^2(1+r-r^2) \approx 0.1159,
  \label{eq:as}
\end{equation}
residual $-1.71\%$ ($0.58\,\varepsilon_{\text{floor}}$) versus the PDG value
$0.1179$.

\textbf{Accumulated braiding floor.} A quantity requiring $n$ recursion cycles
to derive accumulates a precision floor of $n\,\varepsilon_{\text{floor}}$.
This is the correct denominator for normalising residuals of cycle-depth
quantities. A residual within $n\,\varepsilon_{\text{floor}}$ is consistent
with the framework; a residual well below $n\,\varepsilon_{\text{floor}}$ would
falsify it.

%% ────────────────────────────────────────────────────────────
\section{The $\varphi^4$ Colour-Cycle Pattern}
%% ────────────────────────────────────────────────────────────

In UM the charged-lepton hierarchy arises from one recursion cycle per
generation step ($\varphi^1$ spacing). The quark sector traverses three
additional colour-sector channels, giving four cycles per generation step
and a bare mass spacing of $\varphi^4$.

For each quark flavour $q$, define the cycle coordinate
$\ell_q = \log_\varphi(m_q/m_e)$ and the nearest integer multiple of~4,
$n_q = 4\,\text{round}(\ell_q/4)$. The bare UM prediction is
$\varphi^{n_q}$ for the ratio $m_q/m_e$.

\begin{center}
\begin{tabular}{lrrrrrr}
\toprule
Quark & $m_q/m_e$ & $\ell_q$ & $n_q$ &
  $\varphi^{n_q}$ & $|\ell_q - n_q|$ & $n_q\,\varepsilon_\text{floor}$ \\
\midrule
$u$ & 4.2  & 3.00 & 4  & 6.85   & 1.00 & $11.8\%$ \\
$d$ & 9.1  & 4.60 & 4  & 6.85   & 0.60 & $11.8\%$ \\
$s$ & 182.8 & 10.82 & 12 & 322.0 & 1.18 & $35.4\%$ \\
$c$ & 2485 & 16.25 & 16 & 2207  & 0.25 & $47.2\%$ \\
$b$ & 8180 & 18.72 & 20 & 15127 & 1.28 & $59.0\%$ \\
$t$ & 337945 & 26.46 & 28 & 710647 & 1.54 & $82.6\%$ \\
\bottomrule
\end{tabular}
\end{center}

Using PDG 2022 $\overline{\text{MS}}$ masses at their conventional reference
scales ($\mu = 2\,\text{GeV}$ for $u,d,s$; $\mu = m_q$ for $c,b,t$) and
$m_e = 0.511\,\text{MeV}$.

The $\varphi^4$ grid is visible throughout: all six quarks lie within
$1.6$ cycles of the nearest gridpoint. The mass-ratio residuals in the
table below are normalised by the accumulated floor $n_q\,\varepsilon_{\text{floor}}$:

\begin{center}
\begin{tabular}{lrrrr}
\toprule
Quark & Bare residual & $n_q\,\varepsilon_\text{floor}$ &
  Residual / floor & Status \\
\midrule
$u$ & $+62\%$ & $11.8\%$ & 5.3 & EW correction needed \\
$d$ & $-25\%$ & $11.8\%$ & 2.1 & EW correction needed \\
$s$ & $+76\%$ & $35.4\%$ & 2.2 & QCD running needed \\
$c$ & $-11\%$ & $47.2\%$ & 0.24 & \checkmark \\
$b$ & $+85\%$ & $59.0\%$ & 1.44 & \checkmark \\
$t$ & $+110\%$ & $82.6\%$ & 1.33 & \checkmark \\
\bottomrule
\end{tabular}
\end{center}

The heavy quarks ($c$, $b$, $t$) are consistent with the accumulated braiding
floor without any additional input. The light quarks require corrections
discussed below.

%% ────────────────────────────────────────────────────────────
\section{Isospin Splitting of the First-Generation Doublet}
%% ────────────────────────────────────────────────────────────

\subsection{Derivation}

The $u$ and $d$ quarks form the first-generation $\mathrm{SU}(2)_L$ doublet
and share the same bare cycle depth $n = 4$. Their mass difference is
electroweak, not QCD: the $u$--$d$ splitting is generated by the different
weak-isospin assignments within the doublet.

In UM, $\mathrm{SU}(2)_L$ acts through the boundary channel ($W_B$) and
$\mathrm{U}(1)_Y$ acts through the matter channel ($W_M$). The Yukawa coupling
of a quark is proportional to the square root of its channel weight:
\begin{equation}
  y_q \propto \sqrt{W_q},
\end{equation}
because the mass operator requires a single coupling insertion at each vertex.
For the $\mathrm{SU}(2)_L$ doublet components:
\begin{align}
  u\;(T_3 = +\tfrac{1}{2})&:
    \text{does not flip between sectors}
    \;\Rightarrow\; y_u \propto W_M^{1/2}, \\
  d\;(T_3 = -\tfrac{1}{2})&:
    \text{crosses the boundary channel}
    \;\Rightarrow\; y_d \propto W_B^{1/2}.
\end{align}

Since $m_q = y_q\,v/\sqrt{2}$, the ratio is
\begin{equation}
  \boxed{\frac{m_d}{m_u} = \sqrt{\frac{W_B}{W_M}} = \sqrt{\frac{2(1-r)}{r}}.}
  \label{eq:udrat}
\end{equation}

This is equivalent to $1/\!\sin\theta_W$, since the Weinberg angle result
(Paper~4) gives $\sin^2\!\theta_W = W_M/W_B$, hence
$\sqrt{W_B/W_M} = 1/\!\sin\theta_W$.

The expression simplifies to $\sqrt{4\varphi - 2} = \sqrt{2\sqrt{5}}$:
\begin{equation}
  \frac{m_d}{m_u} = \sqrt{2\sqrt{5}} \approx 2.1147.
  \label{eq:udval}
\end{equation}

\subsection{Numerical result}

PDG 2022: $m_d/m_u = 4.67/2.16 = 2.162$.
\[
  \text{Residual} = -2.2\% = 0.74\,\varepsilon_{\text{floor}}.
\]

\subsection{Individual masses}

Both $u$ and $d$ arise from the same bare isospin-symmetric mass
$m_0 = \varphi^4 m_e = 3.502\,\text{MeV}$, which equals the arithmetic
isospin average $(m_u + m_d)/2 = 3.415\,\text{MeV}$ to $2.6\%$
($0.85\,\varepsilon_{\text{floor}}$). The EW splitting distributes this
around the mean:
\begin{equation}
  m_u^{\text{UM}} = \frac{2m_0}{1 + m_d/m_u} = 2.249\,\text{MeV},
  \qquad
  m_d^{\text{UM}} = \frac{2m_0\,(m_d/m_u)}{1 + m_d/m_u} = 4.756\,\text{MeV}.
\end{equation}

\begin{center}
\begin{tabular}{lrrrr}
\toprule
 & UM & PDG & Error & $n\,\varepsilon_\text{floor}$ units \\
\midrule
$m_d/m_u$  & 2.1147  & 2.162  & $-2.2\%$  & 0.74 \\
$m_u$      & 2.249\,MeV & 2.16\,MeV & $+4.1\%$ & 0.35 \\
$m_d$      & 4.756\,MeV & 4.67\,MeV & $+1.8\%$ & 0.16 \\
\bottomrule
\end{tabular}
\end{center}

After applying the EW splitting, both $u$ and $d$ are well within the
4-cycle accumulated floor of $11.8\%$.

%% ────────────────────────────────────────────────────────────
\section{The Strange Quark}
%% ────────────────────────────────────────────────────────────

The strange quark sits at cycle depth $n_s = 12$, four cycles below charm
($n_c = 16$). Unlike the $u$--$d$ pair, $s$ and $c$ are not
degenerate on the bare $\varphi^4$ grid: their 4-cycle separation is a
colour-sector step, not an EW correction.

\subsection{Bare prediction and PDG uncertainty}

The bare $\varphi^{12} m_e = 164.5\,\text{MeV}$ overpredicts the PDG central
value $m_s(2\,\text{GeV}) = 93.4\,\text{MeV}$ by $76\%$, or
$2.15\,n\,\varepsilon_{\text{floor}}$ at $n = 12$. However, the PDG
1$\sigma$ uncertainty is $\pm 11\,\text{MeV}$ (approximately $\pm 12\%$ of the
central value), reflecting the difficulty of extracting $m_s$ via lattice QCD
at scales where $\alpha_s \approx 0.35$. At the $1\sigma$ upper bound of
$104\,\text{MeV}$:
\[
  \text{Residual}_{1\sigma} = \frac{164.5 - 104.4}{104.4} = +57.6\%
  = 1.63\,n\,\varepsilon_{\text{floor}}.
\]
The strange quark is therefore \emph{within} the accumulated floor at the
$1\sigma$ level of experimental uncertainty.

\subsection{Boundary-channel estimate}

The strange quark is the second-generation down-type quark. Like the $d$
quark, it couples to the QCD vacuum through the boundary channel. The
boundary-channel suppression factor is $\sqrt{W_B}$ (the amplitude for a
single boundary-channel coupling), giving:
\begin{equation}
  m_s^{\text{BC}} = \varphi^{12}\,m_e\,\sqrt{W_B}
  = 164.5 \times 0.6535 = 107.5\,\text{MeV}.
  \label{eq:ms_bc}
\end{equation}

\[
  \text{Residual} = \frac{107.5 - 93.4}{93.4} = +15.1\%
  = 0.43\,n\,\varepsilon_{\text{floor}} \quad (n = 12).
\]

This estimate is well within the accumulated floor and is consistent with
the PDG value at the $1\sigma$ level ($m_s = 93.4 \pm 11\,\text{MeV}$:
the $\sqrt{W_B}$ prediction of $107.5\,\text{MeV}$ lies $1.3\sigma$ above
the central value). The $\sqrt{W_B}$ suppression arises from the
boundary-channel coupling of the strange quark to the QCD condensate,
analogous to the $\sqrt{W_B/W_M}$ splitting that governs the $u$--$d$ doublet.

The rigorous derivation of the $\sqrt{W_B}$ prefactor from the $G_2$
zero-mode structure is the next calculation in the strange-quark programme.
This requires the winding-number integral for the strange quark's position
in the $G_2$ cycle hierarchy, which is the subject of the companion
heterotic paper.

%% ────────────────────────────────────────────────────────────
\section{Complete Quark Sector Summary}
%% ────────────────────────────────────────────────────────────

\begin{center}
\begin{tabular}{llrrrr}
\toprule
Quark & UM expression & UM value & PDG & Error & $n\,\varepsilon$ units \\
\midrule
$m_d/m_u$  & $\sqrt{W_B/W_M}$             & 2.115    & 2.162   & $-2.2\%$  & 0.74 \\
$m_u$      & $2\varphi^4 m_e / (1{+}\rho)$  & 2.249\,MeV & 2.16\,MeV & $+4.1\%$ & 0.35 \\
$m_d$      & $2\varphi^4 m_e \rho/(1{+}\rho)$ & 4.756\,MeV & 4.67\,MeV & $+1.8\%$ & 0.16 \\
$m_s$      & $\varphi^{12}m_e\sqrt{W_B}$ (BC)  & 107.5\,MeV & 93.4\,MeV & $+15.1\%$ & 0.43 \\
$m_c/m_e$  & $\varphi^{16}$ (bare)         & 2207     & 2485    & $-11.2\%$ & 0.24 \\
$m_b/m_e$  & $\varphi^{20}$ (bare)         & 15127    & 8180    & $+85\%$   & 1.44 \\
$m_t/m_e$  & $\varphi^{28}$ (bare)         & 710647   & 337945  & $+110\%$  & 1.33 \\
$\alpha_s(M_Z)$ & $r^2(1{+}r{-}r^2)$       & 0.1159   & 0.1179  & $-1.7\%$  & 0.58 \\
\bottomrule
\end{tabular}
\end{center}

where $\rho = \sqrt{W_B/W_M} = \sqrt{2\sqrt{5}}$ and all residuals are
normalised by the accumulated floor $n\,\varepsilon_{\text{floor}}$ at the
relevant cycle depth $n$.

%% ────────────────────────────────────────────────────────────
\section{Discussion}
%% ────────────────────────────────────────────────────────────

\subsection{What is derived}

The $\varphi^4$ cycle structure derives the heavy-quark mass hierarchy
($c$, $b$, $t$) within their accumulated braiding floors without any QCD
input beyond $\alpha_s$. The $u$--$d$ splitting follows from the Weinberg
channel structure already derived in Paper~4; no additional parameters are
introduced.

\subsection{What remains}

The strange quark at $2.1\,n\,\varepsilon_{\text{floor}}$ is the frontier.
Its large relative uncertainty ($\pm 12\%$) means the bare prediction is
consistent with the $2\sigma$ PDG band. The full treatment requires running
from the string scale ($M_\text{Pl}\,\varphi^{-14}$) using the UM
$\alpha_s$ trajectory. This is the next calculation in the quark-mass
programme.

\subsection{Comparison with the Standard Model}

The SM accepts all six quark masses as independent Yukawa couplings —
six free parameters. UM derives $u$ and $d$ individually from two
structural results ($\varphi^4 m_e$ and $\sqrt{W_B/W_M}$), and $c$, $b$, $t$
from the $\varphi^4$ grid alone. The one remaining free assignment is the
integer cycle depth $n_q$ per flavour, which is fixed by the observed
generation ordering and contains no continuous freedom.

\end{document}
